% Options for packages loaded elsewhere
\PassOptionsToPackage{unicode}{hyperref}
\PassOptionsToPackage{hyphens}{url}
\PassOptionsToPackage{dvipsnames,svgnames,x11names}{xcolor}
%
\documentclass[
  11pt,
]{article}
\usepackage{amsmath,amssymb}
\usepackage{iftex}
\ifPDFTeX
  \usepackage[T1]{fontenc}
  \usepackage[utf8]{inputenc}
  \usepackage{textcomp} % provide euro and other symbols
\else % if luatex or xetex
  \usepackage{unicode-math} % this also loads fontspec
  \defaultfontfeatures{Scale=MatchLowercase}
  \defaultfontfeatures[\rmfamily]{Ligatures=TeX,Scale=1}
\fi
\usepackage{lmodern}
\ifPDFTeX\else
  % xetex/luatex font selection
\fi
% Use upquote if available, for straight quotes in verbatim environments
\IfFileExists{upquote.sty}{\usepackage{upquote}}{}
\IfFileExists{microtype.sty}{% use microtype if available
  \usepackage[]{microtype}
  \UseMicrotypeSet[protrusion]{basicmath} % disable protrusion for tt fonts
}{}
\makeatletter
\@ifundefined{KOMAClassName}{% if non-KOMA class
  \IfFileExists{parskip.sty}{%
    \usepackage{parskip}
  }{% else
    \setlength{\parindent}{0pt}
    \setlength{\parskip}{6pt plus 2pt minus 1pt}}
}{% if KOMA class
  \KOMAoptions{parskip=half}}
\makeatother
\usepackage{xcolor}
\usepackage[left=1.25in,right=1.25in,top=1in,bottom=1in]{geometry}
\usepackage{longtable,booktabs,array}
\usepackage{calc} % for calculating minipage widths
% Correct order of tables after \paragraph or \subparagraph
\usepackage{etoolbox}
\makeatletter
\patchcmd\longtable{\par}{\if@noskipsec\mbox{}\fi\par}{}{}
\makeatother
% Allow footnotes in longtable head/foot
\IfFileExists{footnotehyper.sty}{\usepackage{footnotehyper}}{\usepackage{footnote}}
\makesavenoteenv{longtable}
\setlength{\emergencystretch}{3em} % prevent overfull lines
\providecommand{\tightlist}{%
  \setlength{\itemsep}{0pt}\setlength{\parskip}{0pt}}
\setcounter{secnumdepth}{-\maxdimen} % remove section numbering
\ifLuaTeX
  \usepackage{selnolig}  % disable illegal ligatures
\fi
\IfFileExists{bookmark.sty}{\usepackage{bookmark}}{\usepackage{hyperref}}
\IfFileExists{xurl.sty}{\usepackage{xurl}}{} % add URL line breaks if available
\urlstyle{same}
\hypersetup{
  pdftitle={Interim report --- verification and campaign work to date},
  pdfauthor={Trey Morris with Claude Opus 4.7},
  colorlinks=true,
  linkcolor={blue},
  filecolor={Maroon},
  citecolor={Blue},
  urlcolor={blue},
  pdfcreator={LaTeX via pandoc}}

\title{Interim report --- verification and campaign work to date}
\author{Trey Morris with Claude Opus 4.7}
\date{2026-06-01}

\begin{document}
\maketitle

\hypertarget{cover-note}{%
\section{§1 Cover note}\label{cover-note}}

Following our phone call, we consolidate here the verification work and
the three downstream campaigns undertaken since 2026-05-11. The intent
is a document that can be read in one sitting and responded to in
another.

The structure is as follows. §2 describes the process. §3 records the
three Jackson results we believe will interest you most. §4 lists the
two experimental regimes the Jackson work surfaced as candidates for
further attention. §5 collects the broader unconditional and conditional
results across the four campaigns. §6 carries the single most
load-bearing question we have for you, namely the resolution of the
published value of \(r_{e}\). §7 records the secondary errata findings
and citation questions. §8 closes.

A note on numbering. Tags of the form \texttt{\#NN} in the original
draft of this report referred to GitHub issues and pull requests in the
public \texttt{temoTxt/PyPhysics} repository. We have moved every such
tag into a numbered footnote, together with the file paths and external
links that the campaign documents cite. The bibliography section at the
end resolves each footnote.

\hypertarget{papers-and-source-texts-referenced}{%
\subsection{Papers and source texts
referenced}\label{papers-and-source-texts-referenced}}

The table below lists every paper and textbook cited by abbreviation, so
that the abbreviations are unambiguous in what follows. The first group
covers the dual-theory papers from your corpus. The second group covers
other primary sources used in §5 (Bethe 1947 supplies the Lamb-shift
route). The third group covers the campaign textbooks (Jackson,
Griffiths, Bethe-Salpeter).

\begin{longtable}[]{@{}
  >{\raggedright\arraybackslash}p{(\columnwidth - 8\tabcolsep) * \real{0.2000}}
  >{\raggedright\arraybackslash}p{(\columnwidth - 8\tabcolsep) * \real{0.2000}}
  >{\raggedright\arraybackslash}p{(\columnwidth - 8\tabcolsep) * \real{0.2000}}
  >{\raggedright\arraybackslash}p{(\columnwidth - 8\tabcolsep) * \real{0.2000}}
  >{\raggedright\arraybackslash}p{(\columnwidth - 8\tabcolsep) * \real{0.2000}}@{}}
\toprule\noalign{}
\begin{minipage}[b]{\linewidth}\raggedright
Abbreviation
\end{minipage} & \begin{minipage}[b]{\linewidth}\raggedright
Full title
\end{minipage} & \begin{minipage}[b]{\linewidth}\raggedright
Authors
\end{minipage} & \begin{minipage}[b]{\linewidth}\raggedright
Year
\end{minipage} & \begin{minipage}[b]{\linewidth}\raggedright
Venue
\end{minipage} \\
\midrule\noalign{}
\endhead
\bottomrule\noalign{}
\endlastfoot
Maxwell paper & \emph{Two Mathematically Equivalent Versions of
Maxwell's Equations} & Gill, Zachary & 2011 & \emph{Foundations of
Physics} \textbf{41}, 99--128 \\
Analytic Dirac & \emph{Analytic Representation of the Dirac Equation} &
Gill, Zachary & ca. 2002 & \emph{Foundations of Physics Letters}
(subject to author confirmation) \\
TCEP & \emph{The Classical Electron Problem} & Gill, Zachary, Lindesay &
2001 & \emph{Foundations of Physics} \textbf{31}, 1299--1354 \\
DRQM I & \emph{Dual Relativistic Quantum Mechanics I} & Gill, Ares de
Parga, Morris, Wade & 2021 & Manuscript dated 21 August 2021 (Howard
University) \\
Bethe (1947) & \emph{The Electromagnetic Shift of Energy Levels} & Bethe
& 1947 & \emph{Physical Review} \textbf{72}, 339 \\
Jackson & \emph{Classical Electrodynamics} (3rd ed.) & Jackson & 1998 &
Wiley \\
Griffiths 3e & \emph{Introduction to Quantum Mechanics} (3rd ed.) &
Griffiths, Schroeter & 2018 & Cambridge University Press \\
Bethe-Salpeter & \emph{Quantum Mechanics of One- and Two-Electron Atoms}
& Bethe, Salpeter & 1957 / 1977 reprint & Plenum 1957; Springer 1977
reprint \\
\end{longtable}

The Analytic Dirac venue is left tentative because the converted-PDF
metadata does not record the publication line cleanly. One of the
questions in §7 (Q6) is whether the citation we have matches the journal
record you intend.

\hypertarget{how-we-got-here-the-process}{%
\section{§2 How we got here --- the
process}\label{how-we-got-here-the-process}}

The work began in May 2026 with a systematic Wolfram-Mathematica
re-derivation of every numbered equation in the dual-theory paper
corpus. The corpus covers the Maxwell paper, the Analytic Dirac paper,
DRQM I, and TCEP. The Banach-space and Navier-Stokes mathematics papers
were verified at the load-bearing-results level only. The output is a
per-paper verification document\footnote{Per-paper verification
  documents, \texttt{Roadmapping/Equation\_Verification/}, repository
  \texttt{temoTxt/PyPhysics}.} recording each equation as it appears in
the paper, a single-line Wolfram check, a step-by-step expanded
derivation, a cross-reference to the standard textbook equivalent
(Jackson, Sakurai, Goldstein, or Weinberg as applicable), and a verdict.

That verification campaign produced three substantive findings flagged
for author review (§7) and three load-bearing structural results that
survived the re-derivation cleanly. The verification work then opened
three downstream campaigns. The first is the Electromagnetism / Jackson
campaign\footnote{Electromagnetism / Jackson canonical-problems
  campaign. Issue \#42, \texttt{github.com/temoTxt/PyPhysics/issues/42}.},
which works through 66 canonical Jackson problems in CGS and SI under
the proper-time Maxwell formulation, spanning Chapters 1--16. The second
is the Quantum Mechanics / Griffiths campaign\footnote{Quantum Mechanics
  / Griffiths canonical-problems campaign. Issue \#49,
  \texttt{github.com/temoTxt/PyPhysics/issues/49}.}, which works through
41 canonical Griffiths problems across editions 2 and 3 under the
proper-time / dual-theory QM formulation, spanning Chapters 1--12. The
third is the Bethe-Salpeter precision-prediction campaign\footnote{Bethe-Salpeter
  precision-prediction campaign. Issue \#50,
  \texttt{github.com/temoTxt/PyPhysics/issues/50}.}, which is the
precision-experiment counterpart to Griffiths and covers 28 results
across the hydrogen, helium, Lamb-shift, hyperfine, positronium, and
muonium chapters. The three campaigns are recorded in three pull
requests\footnote{Electromagnetism per-PR campaign record. PR \#51,
  \texttt{github.com/temoTxt/PyPhysics/pull/51}.}\footnote{Quantum
  Mechanics per-PR campaign record. PR \#52,
  \texttt{github.com/temoTxt/PyPhysics/pull/52}.}\footnote{Bethe-Salpeter
  per-PR campaign record. PR \#53,
  \texttt{github.com/temoTxt/PyPhysics/pull/53}.}.

Two discipline choices we wish to flag. First, all AI-assisted work in
the repository carries Crocco compliance markers per the
\emph{Responsible AI in Non-Empirical Research} framework\footnote{Crocco,
  Rasdi, and Garavan, \emph{Responsible AI in Non-Empirical Research}
  (2026). Repository compliance mapping:
  \texttt{Roadmapping/Tooling/CROCCO\_COMPLIANCE.md}.}. Under that
framework, substantive AI work (the kind that shapes what a manuscript
argues) requires per-paragraph review blocks and a documented
human-acceptance pass; pragmatic AI work (transcription, Mathematica
verification, formatting) does not. The present report is substantive AI
end-to-end. Second, every per-PR document carries an explicit
honest-framing discipline (the §13 review introduced in the Jackson
plan) that distinguishes algebraic identities from physical claims,
conditional predictions from unconditional ones, and within-scope
verdicts from out-of-scope ones. The Phase 3.5 devil's-advocate pass on
the present report is the same discipline applied to the report itself.

\hypertarget{jackson-highlights-initial-results}{%
\section{§3 Jackson highlights --- initial
results}\label{jackson-highlights-initial-results}}

The Electromagnetism campaign\footnote{Electromagnetism / Jackson
  canonical-problems campaign. Issue \#42,
  \texttt{github.com/temoTxt/PyPhysics/issues/42}.} worked through 66
canonical Jackson problems under the proper-time dual-theory formulation
of the Maxwell paper. Three results merit foregrounding.

\hypertarget{problem-14.2-the-proper-time-liuxe9nard-wiechert-third-term}{%
\subsection{§3.1 Problem 14.2 --- the proper-time Liénard-Wiechert third
term}\label{problem-14.2-the-proper-time-liuxe9nard-wiechert-third-term}}

We observe that Eq. (7) of the Maxwell paper contains a third term in
the radiation-zone fields, \[
\frac{e\,(\mathbf{u}\cdot\mathbf{a})\,[\mathbf{r}\times(\mathbf{u}\times\mathbf{r})]}{b^{4}\,s^{3}},
\qquad b = \sqrt{c^{2} + \mathbf{u}^{2}},
\tag{1}
\] which is absent from the classical Liénard-Wiechert expression. The
coefficient is suppressed by \((\mathbf{u}\cdot\mathbf{a})/b^{4}\) at
non-relativistic intensity. It contributes at order unity for
\(\mathbf{u}/c \approx 1\).

Working through Jackson Problem 14.2 under the dual-theory formulation,
one obtains the explicit angular distribution of the radiation under
(1). It follows that the formulation predicts a \textbf{longitudinal
radiation component} in the far field, which the classical
Liénard-Wiechert theory does not. The classical expression is purely
transverse in the radiation zone.

The result is a structural prediction, not yet an empirical claim.
Indeed, it is (to the extent we have looked) the most natural place to
seek an operational signature of the dual-theory radiation theory. The
Jackson per-PR document under the Electromagnetism PR\footnote{Electromagnetism
  per-PR campaign record. PR \#51,
  \texttt{github.com/temoTxt/PyPhysics/pull/51}.} records the algebra
and the cross-reference to Eq. (7) of the Maxwell paper.

\hypertarget{problem-16.1-abraham-lorentz-dissolution}{%
\subsection{§3.2 Problem 16.1 --- Abraham-Lorentz
dissolution}\label{problem-16.1-abraham-lorentz-dissolution}}

The proper-time formulation supplies the radiation-reaction force as a
first-order \(\partial_{\tau}\) correction. The dissipative coefficient
is the \((\mathbf{u}\cdot\mathbf{a})/b^{4}\) factor of Eq. (4) of the
Maxwell paper. By contrast, the textbook Abraham-Lorentz force is
third-order in \(\partial_{t}\), \[
\mathbf{F}_{\rm AL} \;=\; \tfrac{2}{3}\,\frac{e^{2}}{c^{3}}\,\dddot{\mathbf{x}}.
\tag{2}
\] The third-order time derivative in (2) is the source of the
runaway-solution and pre-acceleration pathologies that mar the classical
theory.

We observe that the first-order proper-time structure removes both
pathologies. The equation of motion is first-order in \(\tau\) and
therefore does not admit the higher-order spurious solutions. The
Jackson Problem 16.1 derivation makes the dissolution claim explicit.

The claim is structural. Whether it also resolves the empirical defects
that motivated the Cole, Poder, and Wistisen 2018 experiments is what we
believe deserves a separate look. The candidate experiments are recorded
in §4.

\hypertarget{problem-5.13-a-non-zero-result-we-did-not-anticipate}{%
\subsection{§3.3 Problem 5.13 --- a non-zero result we did not
anticipate}\label{problem-5.13-a-non-zero-result-we-did-not-anticipate}}

Working through Jackson's magnetic-moment-of-a-current-loop problem in
the proper-time formulation gave a non-zero result, where the classical
calculation gives zero. The result is flagged with a tagged comment on
the Electromagnetism issue\footnote{Cross-verification tag
  (\texttt{@Zorns-Lemmon}) on the Electromagnetism master issue. Issue
  \#42, \texttt{github.com/temoTxt/PyPhysics/issues/42}.} for
cross-verification. It is not yet checked by a second pair of eyes. The
campaign's per-PR document records the calculation in full so that the
result is reproducible.

\hypertarget{candidate-experiments-what-we-found-to-look-at}{%
\section{§4 Candidate experiments --- what we found to look
at}\label{candidate-experiments-what-we-found-to-look-at}}

The Jackson work surfaced two experimental regimes in which the
dual-theory predictions appear to be operationally distinguishable from
textbook radiation theory at current measurement precision. We have
scoped both as sibling issues to the Electromagnetism master, but have
not yet executed the comparison. We flag them here for your awareness.

\hypertarget{cole-2018-poder-2018-wistisen-2018}{%
\subsection{§4.1 Cole 2018 / Poder 2018 / Wistisen
2018}\label{cole-2018-poder-2018-wistisen-2018}}

Three radiation-reaction experiments operate in the GeV and
extreme-intensity regime. Cole and Poder use the Gemini laser-wakefield
accelerator with a counter-propagating laser pulse at intensity \[
I \;\sim\; 10^{21}\;{\rm W/cm^{2}}.
\tag{3}
\] Wistisen uses the CERN SPS 200 GeV electron beam channeling through
aligned silicon. The three papers themselves disagree on which classical
limit fits their data: Cole and Poder favour quantum-corrected
Landau-Lifshitz, while Wistisen prefers classical Landau-Lifshitz. The
disagreement is itself a signal that the regime is ambiguous. The
proper-time third term of (1) and the dissipative coefficient of (2)
make distinguishable predictions in both geometries. The scoping is
recorded on the dedicated sibling issue\footnote{Cole 2018 / Poder 2018
  / Wistisen 2018 radiation-reaction scoping. Issue \#43,
  \texttt{github.com/temoTxt/PyPhysics/issues/43}.}.

\hypertarget{mev-bremsstrahlung-at-clinical-linac-energies}{%
\subsection{§4.2 MeV bremsstrahlung at clinical-linac
energies}\label{mev-bremsstrahlung-at-clinical-linac-energies}}

A precision-experiment regime distinct from §4.1 emerges from
medical-physics calibration. Calibration measurements at clinical-linac
electron energies (typically 1, 6, and 18 MeV) deliver
\(\sim 1\)--\(2\%\) precision on the bremsstrahlung spectrum across
decades of data, against the NIST EGSnrc and AAPM Task Group standards.
At MeV electron energies the third term of (1) contributes at order
unity, with \[
\frac{u}{c} \;\approx\; 0.94 \quad\text{at}\quad E_{e} = 1\;{\rm MeV}.
\tag{4}
\] The prediction is a longitudinal polarisation component absent from
classical Born-approximation QED. The scoping is recorded on the
dedicated sibling issue\footnote{MeV-bremsstrahlung clinical-linac
  scoping. Issue \#48, \texttt{github.com/temoTxt/PyPhysics/issues/48}.}.
The bibliography stubs for the target references (Seltzer-Berger tables,
Pratt-Tseng review, NIST EGSnrc) are not yet in place.

\hypertarget{what-the-campaigns-found-unconditional-and-conditional}{%
\section{§5 What the campaigns found --- unconditional and
conditional}\label{what-the-campaigns-found-unconditional-and-conditional}}

Three structural results across the four campaigns are unconditional, in
the sense that they do not depend on the resolution of the \(r_{e}\)
branch question of §6.

First, the algebraic structure of the Maxwell paper (Eqs. (22)--(24),
modulo the typos recorded under Finding 1 of §7) reproduces under
Wolfram verification.

Second, the §II Foldy-Wouthuysen reduction of the dual-Dirac equation in
DRQM I produces the standard Pauli Hamiltonian at leading
\((Z\alpha)^{2}\) order, with the standard spin-orbit,
relativistic-kinetic, and Darwin coefficients.

Third, Bethe's 1947 mass-renormalisation estimate of the
2S\(_{1/2}\)--2P\(_{1/2}\) Lamb shift reproduces under the proper-time
matrix elements. We obtain the campaign's predicted Lamb shift, \[
\Delta\nu_{\rm Lamb}^{\rm framework}
\;\approx\; 1{,}016\;{\rm MHz},
\tag{5}
\] which is the same number that the standard Bethe-estimate route
gives. The residual to the measured value, \[
\Delta\nu_{\rm Lamb}^{\rm meas} \;=\; 1{,}057.845(9)\;{\rm MHz},
\tag{6}
\] is \(\sim 42\) MHz. The residual is attributable to the sub-leading
one-loop QED contributions dropped by the Bethe estimate, which are out
of scope for the campaign's apparatus.

The remaining results across the Bethe-Salpeter campaign\footnote{Bethe-Salpeter
  precision-prediction campaign. Issue \#50,
  \texttt{github.com/temoTxt/PyPhysics/issues/50}.} are
\emph{conditional} on the resolution of the \(r_{e}\) finding. Six
independent precision atomic-physics observables exhibit a consistent
branched verdict, set out in the table below.

\begin{longtable}[]{@{}
  >{\raggedright\arraybackslash}p{(\columnwidth - 8\tabcolsep) * \real{0.2000}}
  >{\raggedright\arraybackslash}p{(\columnwidth - 8\tabcolsep) * \real{0.2000}}
  >{\raggedright\arraybackslash}p{(\columnwidth - 8\tabcolsep) * \real{0.2000}}
  >{\raggedright\arraybackslash}p{(\columnwidth - 8\tabcolsep) * \real{0.2000}}
  >{\raggedright\arraybackslash}p{(\columnwidth - 8\tabcolsep) * \real{0.2000}}@{}}
\toprule\noalign{}
\begin{minipage}[b]{\linewidth}\raggedright
Observable
\end{minipage} & \begin{minipage}[b]{\linewidth}\raggedright
Branch (b) prediction
\end{minipage} & \begin{minipage}[b]{\linewidth}\raggedright
Branch (c) prediction
\end{minipage} & \begin{minipage}[b]{\linewidth}\raggedright
Measurement
\end{minipage} & \begin{minipage}[b]{\linewidth}\raggedright
Uncertainty
\end{minipage} \\
\midrule\noalign{}
\endhead
\bottomrule\noalign{}
\endlastfoot
Electron \(g_{s}\) & \(-2.000\,571\,4\) & \(-2.002\,319\,30\) &
\(-2.002\,319\,30\ldots\) & \(\sim 10^{-12}\) \\
H \(2P_{3/2} - 2P_{1/2}\) & \(\sim 10{,}952\) MHz & \(\sim 10{,}962\)
MHz & \(10{,}969.13(10)\) MHz & \(\sim 10^{-8}\) \\
H 1S hyperfine (21 cm) & \(\sim 1{,}418.81\) MHz & \(\sim 1{,}420.04\)
MHz & \(1{,}420.405\,751\,768(2)\) MHz & \(\sim 10^{-12}\) \\
He \({}^{3}P_{0} - {}^{3}P_{1}\) & \(\sim 29{,}617.4\) MHz &
\(\sim 29{,}616.95\) MHz & \(29{,}616.952\) MHz & \(\sim 10^{-9}\) \\
Positronium ortho-para & \(\sim 203{,}505\) MHz & \(\sim 203{,}389\) MHz
& \(203{,}389(2)\) MHz & \(\sim 10^{-5}\) \\
Muonium hyperfine & \(\sim 4{,}464.6\) MHz & \(\sim 4{,}463.4\) MHz &
\(4{,}463.302\,776(51)\) MHz & \(\sim 10^{-8}\) \\
\end{longtable}

Across all six observables, branch (b) is in \(\sim 10^{-3}\) fractional
disagreement with measurement, while branch (c) agrees with measurement
at the campaign's Bethe-estimate precision floor. The pattern is the
consequence of the \(r_{e}\) finding propagating through every
\(g_{s}\)-linear and \(g_{s}^{2}\)-linear observable in the same way.

The Bethe-estimate Lamb shift of (5)--(6) is the campaign's clearest
measurement-level non-falsification at the precision the route can
deliver. The framework reproduces the same \(\sim 1{,}016\) MHz number
that the standard QED Bethe-estimate route gives, and the \(\sim 42\)
MHz residual is shared between the two routes (it is what neither
captures without a full one-loop calculation). It is not, in the strict
sense, an endorsement that distinguishes the framework from standard
QED. The two formulations give the same Bethe-estimate number, and the
campaign's apparatus cannot push below this precision floor. A
proper-time one-loop dual-QED calculation is documented as future
work\footnote{Proper-time one-loop dual-QED Lamb-shift calculation
  (future work). Issue \#55,
  \texttt{github.com/temoTxt/PyPhysics/issues/55}.}. The relevant point
is that the framework does not fail at the Bethe-estimate precision.
That is the strongest empirical claim the present apparatus supports.
The six observables in the table, by contrast, do depend on \(r_{e}\) at
leading order.

\hypertarget{the-r_e-branch-question-the-load-bearing-ask}{%
\section{\texorpdfstring{§6 The \(r_{e}\) branch question --- the
load-bearing
ask}{§6 The r\_\{e\} branch question --- the load-bearing ask}}\label{the-r_e-branch-question-the-load-bearing-ask}}

This is the single most consequential question we have for you. We have
carried two values of \(r_{e}/r_{0}\) through the campaigns, because the
published value and the value required to reproduce the measured
\(g_{e}\) disagree at the fourth decimal place.

\hypertarget{what-the-two-branches-are}{%
\subsection{§6.1 What the two branches
are}\label{what-the-two-branches-are}}

DRQM I §III.D records the anomalous-\(g\)-factor formula \[
g_{r} \;=\; 2\!\left(1 - \frac{4\,r_{0}}{2r + r_{0}}\right),
\tag{7}
\] together with a stated value \[
\left(\frac{r_{e}}{r_{0}}\right)_{\!\rm published} \;=\; 0.499\,857\,150\,068\,631,
\tag{8}
\] that the paper claims yields \[
g_{e}^{\rm claimed} \;=\; -2.002\,319\,304\,362\,56,
\tag{9}
\] which is the experimental value of \(g_{e}\). We have evaluated (7)
at the value (8) in Wolfram and instead obtain \[
g_{e}^{\rm Wolfram} \;=\; -2.000\,571\,481\,361\,548\,7.
\tag{10}
\] The value of \(r_{e}/r_{0}\) that does reproduce (9) is \[
\left(\frac{r_{e}}{r_{0}}\right)_{\!\rm corrected} \;\approx\; 0.499\,420\,509\,912\,831\,8.
\tag{11}
\] The two values (8) and (11) are very close --- they differ in the
fourth decimal place --- but the formula (7) has a sensitivity \[
\left.\frac{d g_{r}}{d(r_{e}/r_{0})}\right|_{r_{e}/r_{0}\approx 1/2} \;\approx\; 4,
\tag{12}
\] so that even small differences in \(r_{e}\) correspond to substantial
differences in \(g_{e}\). We label the as-published value (8)
\textbf{branch (b)}, and the corrected value (11) that reproduces the
measured \(g_{e}\) \textbf{branch (c)}.

\hypertarget{why-the-question-is-load-bearing}{%
\subsection{§6.2 Why the question is
load-bearing}\label{why-the-question-is-load-bearing}}

The value of \(r_{e}\) enters every observable that depends on the
electron's anomalous magnetic moment. That is every spin-dependent
precision observable in atomic and molecular spectroscopy. The
six-observable table in §5 shows the pattern. Branch (b) consistently
predicts a \(\sim 10^{-3}\) disagreement; branch (c) consistently agrees
with measurement at the Bethe-estimate precision floor.

On the hydrogen 1S hyperfine line --- the 21-cm transition, the most
precisely measured atomic-physics frequency at \(\sim 10^{-12}\)
relative uncertainty --- branch (b)'s disagreement is of order
\(\sim 10^{6}\) standard deviations beyond the measurement uncertainty.

We are not asking you to pick the value that best fits the data. The
campaign's verdict-recording is not a fitting exercise. We are asking
which value the framework intends, so that the campaign can record
verdicts as \textbf{Pass} or \textbf{Refuted} unconditionally rather
than carry both branches indefinitely.

\hypertarget{what-we-are-asking-concretely}{%
\subsection{§6.3 What we are asking,
concretely}\label{what-we-are-asking-concretely}}

The question decomposes into three sub-questions, each of which has a
clear consequence for the campaign.

\textbf{Q1a.} Is the as-published value (8) the value the framework
intends, i.e.~branch (b)? If yes, the campaign's verdicts on the six
observables of §5 would be recorded as \textbf{Refuted} (ruled out at
current measurement precision). The open question would then become how
the framework reconciles the \(\sim 10^{-3}\) disagreement across
hyperfine, fine structure, and \(g\) --- either through a feature of the
derivation we have not captured, or through a regime distinction we have
missed.

\textbf{Q1b.} Is the corrected value (11) the value the framework
intends, i.e.~branch (c)? If yes, the published number (8) in DRQM I
§III.D is a transcription error to be erratum'd, and the campaign's
verdicts collapse to \textbf{Pass} at the Bethe-estimate precision
floor.

\textbf{Q1c.} Is there a third possibility we have not considered --- a
different formula, a different derivation, or a context-dependent choice
where one value is correct in one regime and another in another? If yes,
we will rework the campaign's verdicts with the third value in mind.

\hypertarget{how-the-framework-supports-answering-this}{%
\subsection{§6.4 How the framework supports answering
this}\label{how-the-framework-supports-answering-this}}

Two routes are available. The first is to retrieve the original working
notebook in which \(r_{e}\) was computed and confirm the digit string.
The campaign does not have access to the working notebooks; you do. Even
a one-line confirmation that the published or the corrected value is the
intended one is sufficient for the campaign to move forward.

The second route is to rederive \(r_{e}\) from the renormalisation
prescription of the dual-Dirac equation. This requires the prescription
itself to be documented (the operator-ordering choice in DRQM I §II.D,
which is also Q5 of §7). The campaign could in principle perform the
rederivation under the Weyl ordering we have assumed, but the result is
only as good as the assumption about the prescription.

\hypertarget{secondary-findings-and-questions}{%
\section{§7 Secondary findings and
questions}\label{secondary-findings-and-questions}}

Three additional errata findings and three supporting questions, all of
which are pragmatic-level confirmations rather than load-bearing
structural decisions.

\textbf{Finding 1 --- Maxwell paper Eq. (24): two typos.} The published
Eq. (24) is missing a factor of \(c\) in the \(\Sigma\cdot\mathbf{B}\)
term. The term should read \[
-\,\frac{e\,\hbar\,\Sigma\cdot\mathbf{B}}{2 m c},
\tag{13}
\] which is the value already in Eq. (22). Furthermore, the published
Eq. (24) is missing the cross-term \[
+\,\frac{V^{2}}{2 m c^{2}},
\tag{14}
\] which arises in computing
\(H^{2} \to K = H^{2}/(2 m c^{2}) + m c^{2}/2\). The campaign's Wolfram
check is reproducible, and the algebraic argument is straightforward.

\textbf{Finding 3 --- TCEP Eq. (4.16): sign typo.} Equation (4.16) is
written as \[
v_{g} \;=\; v_{g}' - v,
\tag{15}
\] but the algebra from the paper's own (4.14)--(4.15) gives \[
v_{g} \;=\; v_{g}' + v.
\tag{16}
\] The paper's commentary surrounding the equation also reads as the
\(+\) sign. The sign in (15) appears to be a typesetting error.

\textbf{Q4 --- \(r_{\mu}\) and \(r_{p}\) numerical values for DRQM I Eq.
(III.23).} The formulas for the muon and proton anomalous moments, \[
g_{\mu}^{a}, \qquad g_{p}^{a},
\tag{17}
\] are stated in DRQM I §III.D alongside the electron formula. The
corresponding cutoff values \(r_{\mu}\) and \(r_{p}\) are not given
numerically. Without them, the muonium and proton magnetic-moment
predictions cannot be computed from the framework. If you have these
values, they would close out the analogous discriminator analyses we are
presently unable to do for the muon and proton.

\textbf{Q5 (optional) --- operator-ordering choice in DRQM I §II.D.} The
sub-leading \((Z\alpha)^{4}\) results in the dual-Dirac Foldy-Wouthuysen
reduction depend on the operator-ordering choice. We have assumed Weyl
ordering throughout the campaign, which reproduces the standard Pauli
Hamiltonian at leading order. If the framework intends a different
ordering, the sub-leading results would shift, and the §6.4 route-(ii)
rederivation of \(r_{e}\) would proceed under the different
prescription.

\textbf{Q6 (citation hygiene) --- Analytic Dirac citation details.} The
Analytic Dirac paper's publication venue is recorded in the table at the
head of this report as \emph{Foundations of Physics Letters}, ca. 2002.
The converted-PDF metadata does not record the publication line cleanly.
The cleanest fix would be a one-line confirmation of the journal,
volume, and year as you intend the citation to read, so that the table
can be locked in and any future bibliographic export from the repository
carries the correct line.

\hypertarget{closing}{%
\section{§8 Closing}\label{closing}}

Thank you for the phone call and for the opening to send this. The
single most useful response would be a resolution of §6 (Q1a, Q1b, Q1c).
Even a one-line answer determines which of the campaign's six precision
verdicts collapse to \textbf{Pass} and which become \textbf{Refuted}.
Confirmations or corrections on Finding 1, Finding 3, and Q4
(\(r_{\mu}\), \(r_{p}\)) would close out the verification-level items.
The campaigns continue under the current trajectory if you find the work
useful in this form. If you have a different direction you would prefer,
we will redirect.

\hypertarget{references}{%
\subsection{References}\label{references}}

\end{document}
